% !TeX encoding = UTF-8
\documentclass[institutional,screen,none]{ua-engineering-v6-beamer}
% !TeX encoding = UTF-8
\UASetUniversity{Universidad Austral}
\UASetFaculty{Facultad de Ingeniería}
\UASetDepartment{Departamento de Computación}
\UASetCampus{Pilar}
\UASetCareer{Ingeniería Informática}
\UASetCommission{Comisión A}
\UASetProfessor{Equipo docente}
\UASetSemester{Segundo cuatrimestre 2026}
\UASetCourse{Sistemas Distribuidos}
\UASetDate{2026}

\UASetSemester{Segundo cuatrimestre 2026}
\UASetDate{2026}
\usetikzlibrary{decorations.pathreplacing,matrix,backgrounds,shapes.multipart,fit,positioning,arrows.meta,calc}

\hyphenpenalty=9000
\exhyphenpenalty=9000
\tolerance=1800
\emergencystretch=1.3em

\newcommand{\UASlideTag}[2]{%
  \begin{tikzpicture}[remember picture,overlay]
    \node[anchor=north east,rounded corners=1mm,inner xsep=3pt,inner ysep=1.5pt,
      draw=#2,fill=#2!8,text=#2,font=\fontsize{5.8}{6.4}\selectfont\bfseries]
      at ([xshift=-0.72cm,yshift=-1.08cm]current page.north east) {#1};
  \end{tikzpicture}%
}
\newcommand{\FrameID}[1]{\note{Frame ID: #1}}
\newcommand{\SourceLine}[1]{\vfill{\tiny\color{UAIngMuted}Fuente / referencia: #1}}

\UASetClassNumber{08}
\UASetClassTitle{Causalidad, consenso y consistencia distribuida}
\title{\UACourse}
\subtitle{Clase 08 --- De Kafka y causalidad a consenso, CAP, Dynamo y Spanner}
\author{\UAProfessor}
\date{\UADate}

\tikzset{
  uaBox/.style={draw=UAIngNavy,rounded corners=1.5mm,thick,fill=white,
    minimum height=0.70cm,align=center,font=\scriptsize,inner sep=4pt},
  uaBoxO/.style={uaBox,draw=UAIngOrange,fill=UAIngPaper},
  % Alias de compatibilidad usados por los diagramas de la clase.
  uaOrange/.style={uaBoxO},
  uaSoft/.style={uaBox,draw=UAIngRule,fill=UAIngPaper},
  uaArrow/.style={->,very thick,>=Latex,draw=UAIngNavy},
  uaArrowO/.style={->,very thick,>=Latex,draw=UAIngOrange},
  uaOrangeArrow/.style={uaArrowO},
  uaDashed/.style={->,thick,dashed,>=Latex,draw=UAIngMuted},
  uaLane/.style={draw=UAIngRule,line width=0.9pt},
  uaEvent/.style={circle,fill=UAIngNavy,inner sep=2.2pt},
  uaEventO/.style={circle,fill=UAIngOrange,inner sep=2.4pt},
  uaDot/.style={circle,fill=UAIngNavy,inner sep=2.2pt},
  uaTiny/.style={font=\tiny,align=center,text=UAIngText},
  uaSmall/.style={font=\scriptsize,align=center,text=UAIngText},
  uaState/.style={draw=UAIngNavy,rounded corners=1mm,fill=white,
    minimum width=1.45cm,minimum height=0.48cm,font=\tiny,align=center},
  uaStateO/.style={uaState,draw=UAIngOrange,fill=UAIngPaper}
}

\begin{document}

{
\setbeamertemplate{footline}{}
\begin{frame}[plain]
\titlepage
\end{frame}
}

% ID: C08-00
\begin{frame}{La historia que guiará toda la clase}
\UASlideTag{NÚCLEO}{UAIngNavy}
\scriptsize
\begin{uakeybox}
En esta clase seguiremos el destino de una única entrada VIP que Sofía y Tomás intentan comprar casi al mismo tiempo desde regiones diferentes.
\end{uakeybox}
\vspace{0.15cm}
\begin{columns}[T]
\column{0.48\textwidth}
\begin{block}{Lo que necesitamos comprender}
\begin{itemize}
\item qué eventos están relacionados y cuáles son concurrentes;
\item cómo varias réplicas acuerdan una decisión;
\item qué ocurre cuando la red impide coordinar;
\item qué costo paga cada garantía.
\end{itemize}
\end{block}
\column{0.48\textwidth}
\begin{block}{La pregunta transversal}
¿Qué error es inaceptable y qué costo aceptamos para evitarlo?
\end{block}
\begin{center}
\textbf{latencia · indisponibilidad · coordinación · reparación}
\end{center}
\end{columns}
\end{frame}

% ID: C08-00B
\begin{frame}{Evidencias de aprendizaje}
\UASlideTag{NÚCLEO}{UAIngNavy}
\scriptsize
\begin{columns}[T]
\column{0.48\textwidth}
\begin{block}{Durante la clase}
\begin{itemize}
\item construir \textit{happened-before} y reconocer concurrencia;
\item completar relojes de Lamport y vectoriales;
\item seguir una entrada Raft y predecir un crash o una partición.
\end{itemize}
\end{block}
\column{0.48\textwidth}
\begin{block}{Al cerrar}
\begin{itemize}
\item distinguir mayoría de consenso de quorum Dynamo-like;
\item decidir cuándo un conflicto puede fusionarse;
\item explicar TrueTime y defender una política con falla, costo y métrica.
\end{itemize}
\end{block}
\end{columns}
\begin{uakeybox}
La evidencia será reconstruir una ejecución y justificar por qué el sistema confirma, espera, rechaza, conserva versiones o repara.
\end{uakeybox}
\end{frame}

\UASectionSlide{Puente desde Kafka}

% ID: C08-KF01
\begin{frame}{Repaso visual: un Mundial en tiempo real}
\UASlideTag{NÚCLEO}{UAIngNavy}
\small
\begin{center}
\begin{tikzpicture}[node distance=0.55cm and 0.60cm]
\node[uaOrange,text width=2.0cm] (match) {Partido A\\eventos};
\node[uaSoft,text width=1.8cm,right=of match] (prod) {Productor};
\node[uaSoft,text width=2.0cm,right=of prod] (kafka) {Kafka\\log preservado};
\node[uaSoft,text width=1.8cm,right=of kafka] (cons) {Consumidor};
\node[uaOrange,text width=1.8cm,right=of cons] (score) {Marcador\\web};
\draw[uaOrangeArrow] (match)--node[above,uaTiny]{GoalScored}(prod);
\draw[uaArrow] (prod)--(kafka);
\draw[uaArrow] (kafka)--node[above,uaTiny]{poll}(cons);
\draw[uaOrangeArrow] (cons)--(score);
\end{tikzpicture}
\end{center}
\begin{block}{Idea recuperada}
Kafka preserva una secuencia de registros para que productores y consumidores avancen de manera desacoplada y recuperable.
\end{block}
\end{frame}

% ID: C08-KF02
\begin{frame}{Una única secuencia no alcanza para muchos partidos}
\UASlideTag{NÚCLEO}{UAIngNavy}
\small
\begin{columns}[T]
\column{0.48\textwidth}
\begin{alertblock}{Problema de capacidad}
Muchos partidos compiten por la CPU, red y almacenamiento de una única secuencia.
\end{alertblock}
\column{0.48\textwidth}
\begin{alertblock}{Problema de orden}
Distribuir eventos al azar puede mostrar un gol antes de que comience su partido.
\end{alertblock}
\end{columns}
\vspace{0.3cm}
\begin{center}
\begin{tikzpicture}[node distance=0.42cm]
\node[uaSoft,text width=1.45cm] (a) {Partido A};
\node[uaSoft,text width=1.45cm,below=of a] (b) {Partido B};
\node[uaSoft,text width=1.45cm,below=of b] (c) {Partido C};
\node[uaOrange,text width=3.5cm,right=1.1cm of b] (q) {una sola secuencia\\y un consumidor};
\draw[uaArrow] (a.east)--(q.west);
\draw[uaArrow] (b.east)--(q.west);
\draw[uaArrow] (c.east)--(q.west);
\end{tikzpicture}
\end{center}
\end{frame}

% ID: C08-KF03
\begin{frame}{La key define el ámbito de orden}
\UASlideTag{NÚCLEO}{UAIngNavy}
\scriptsize
\begin{block}{Decisión}
Usamos \texttt{game\_id} como key: todos los eventos de un partido permanecen en la misma partición.
\end{block}
\begin{center}
\begin{tikzpicture}[node distance=0.38cm]
\node[uaOrange,text width=1.25cm] (p0) {P0};
\node[uaSoft,text width=6.2cm,right=of p0] {A: Started $\rightarrow$ YellowCard};
\node[uaOrange,text width=1.25cm,below=0.40cm of p0] (p1) {P1};
\node[uaSoft,text width=6.2cm,right=of p1] {B: Started $\rightarrow$ Substitution };
\node[uaOrange,text width=1.25cm,below=0.40cm of p1] (p2) {P2};
\node[uaSoft,text width=6.2cm,right=of p2] {C: Started $\rightarrow$ Goal};
\end{tikzpicture}
\end{center}
\begin{uakeybox}
Ganamos orden local por partido y paralelismo entre partidos. No obtenemos un offset global gratuito entre P0, P1 y P2.
\end{uakeybox}
\end{frame}

% ID: C08-KF04
\begin{frame}{Consumer groups: repartir trabajo y releer la misma historia}
\UASlideTag{NÚCLEO}{UAIngNavy}
\small
\begin{columns}[T]
\column{0.48\textwidth}
\begin{block}{Grupo \texttt{live-score}}
\begin{itemize}
\item P0 $\rightarrow$ Score-1
\item P1 $\rightarrow$ Score-2
\item P2 $\rightarrow$ Score-3
\end{itemize}
Las particiones se reparten dentro del grupo.
\end{block}
\column{0.48\textwidth}
\begin{block}{Grupo \texttt{analytics}}
Lee nuevamente P0, P1 y P2, pero conserva sus propios offsets.
\end{block}
\end{columns}
\begin{uakeybox}
Una partición tiene una sola instancia activa dentro de un grupo; grupos diferentes pueden leer la misma partición con progreso independiente.
\end{uakeybox}
\end{frame}

% ID: C08-KF05
\begin{frame}{Crash, replay e idempotencia}
\UASlideTag{NÚCLEO}{UAIngNavy}
\small
\begin{center}
\begin{tikzpicture}[node distance=0.48cm]
\node[uaSoft,text width=1.55cm] (poll) {poll\\P2:420};
\node[uaSoft,text width=1.75cm,right=of poll] (effect) {actualiza\\marcador};
\node[uaOrange,text width=1.45cm,right=of effect] (crash) {CRASH};
\node[uaSoft,text width=1.70cm,right=of crash] (replay) {reinicia\\desde 420};
\node[uaSoft,text width=1.65cm,right=of replay] (dedup) {deduplica\\event\_id};
\draw[uaArrow] (poll)--(effect);
\draw[uaOrangeArrow] (effect)--(crash);
\draw[uaArrow] (crash)--(replay);
\draw[uaArrow] (replay)--(dedup);
\end{tikzpicture}
\end{center}
\begin{alertblock}{Límite}
At-least-once permite que el registro reaparezca; no garantiza por sí solo que el efecto lógico ocurra una única vez.
\end{alertblock}
\end{frame}

% ID: C08-KF06
\begin{frame}{Algunas métricas a considerar en Kafka}
\UASlideTag{NÚCLEO}{UAIngNavy}
\scriptsize
\begin{columns}[T]
\column{0.52\textwidth}
\begin{block}{Definiciones}
\begin{itemize}
\item \textbf{RF=3}: líder + dos seguidores configurados.
\item \textbf{ISR}: réplicas suficientemente sincronizadas.
\item \textbf{minISR=2}: hacen falta al menos dos réplicas en ISR.
\item \textbf{acks=all}: todas las réplicas del ISR actual confirman el batch.
\end{itemize}
\end{block}
\begin{uainfo}
Si Réplica3 sale del ISR, pero Réplica1 y Réplica2 siguen sincronizadas, la escritura continúa. Si solo queda Réplica1, se rechaza: no se reconoce éxito con una sola copia sincronizada.
\end{uainfo}
\column{0.45\textwidth}
\begin{tabular}{p{0.50\textwidth}p{0.38\textwidth}}
\toprule
\textbf{Guardrail docente} & \textbf{Referencia}\\
\midrule
mensaje habitual & $<1$ MB\\
datos retenidos/broker & $\sim1$ TB\\
ingreso conservador/broker & $<10$ MB/s\\
\bottomrule
\end{tabular}
\vspace{0.15cm}
\begin{uawarn}
Son presupuestos docentes, no máximos universales. Deben validarse con hardware, tamaño de registros, replicación y pruebas de carga.
\end{uawarn}
\end{columns}
\end{frame}

% ID: C08-KF07
\begin{frame}{Puente: Kafka ordena registros; no elige por sí solo una decisión autoritativa}
\UASlideTag{NÚCLEO}{UAIngNavy}
\scriptsize
\begin{columns}[T]
\column{0.48\textwidth}
\begin{block}{Kafka aporta}
\begin{itemize}
\item orden dentro de cada partición;
\item replay y progreso por grupo;
\item replicación líder--seguidor;
\item desacoplamiento productor--consumidor.
\end{itemize}
\end{block}
\column{0.48\textwidth}
\begin{alertblock}{La nueva pregunta}
Si dos regiones intentan reservar la única entrada VIP, ¿qué comando debe convertirse en la decisión autoritativa?
\end{alertblock}
\end{columns}
\begin{uakeybox}
Pasamos de preservar una historia a acordar qué historia será válida.
\end{uakeybox}
\end{frame}

\UASectionSlide{Por qué estos conceptos se usan hoy}

% ID: C08-R01
\begin{frame}{Sistema real 1: Kubernetes y etcd}
\UASlideTag{CASO REAL}{UAIngOrange}
\small
\begin{block}{Qué es Kubernetes}
Plataforma que administra aplicaciones en contenedores y mantiene un estado deseado del cluster.
\end{block}
\begin{block}{Qué es etcd}
Almacén clave--valor distribuido donde el control plane conserva objetos y configuración autoritativa.
\end{block}
\begin{uakeybox}
etcd utiliza Raft para acordar un log; sin mayoría no puede confirmar nuevos cambios del control plane.
\end{uakeybox}
\end{frame}

% ID: C08-R02
\begin{frame}{Sistema real 2: Amazon Dynamo}
\UASlideTag{CASO REAL}{UAIngOrange}
\small
\begin{block}{Qué fue Dynamo}
Almacén clave--valor altamente disponible diseñado por Amazon para servicios que debían continuar respondiendo durante fallas.
\end{block}
\begin{block}{Decisión de diseño}
Conservar versiones concurrentes, responder con quórums flexibles y reparar después, en vez de bloquear siempre por coordinación global.
\end{block}
\begin{uakeybox}
La disponibilidad se paga con versiones, reconciliación, hints y trabajo de reparación.
\end{uakeybox}
\end{frame}

% ID: C08-R03
\begin{frame}{Sistema real 3: Google Spanner}
\UASlideTag{CASO REAL}{UAIngOrange}
\small
\begin{block}{Qué es Spanner}
Base de datos distribuida globalmente que ofrece transacciones entre grupos de datos y un orden compatible con lo observado por los clientes.
\end{block}
\begin{block}{Decisión de diseño}
Combinar consenso localizado, coordinación transaccional y TrueTime, una abstracción que expone incertidumbre temporal.
\end{block}
\begin{uakeybox}
Spanner paga coordinación y espera temporal para ofrecer garantías globales más fuertes.
\end{uakeybox}
\end{frame}

% ID: C08-R04
\begin{frame}{Lo que un ingeniero decide realmente}
\UASlideTag{NÚCLEO}{UAIngNavy}
\scriptsize
\begin{tabular}{p{0.18\textwidth}p{0.22\textwidth}p{0.22\textwidth}p{0.18\textwidth}}
\toprule
\textbf{Operación} & \textbf{Invariante} & \textbf{Error tolerable} & \textbf{Evidencia}\\
\midrule
última entrada & confirmaciones $\le 1$ & rechazo temporal & doble venta, p99\\
carrito & no perder intención & ramas fusionables & conflictos, convergencia\\
ranking & atraso acotado & vista anterior & lag temporal\\
ledger (registros contables) & no crear dinero & espera o rechazo & anomalías, abortos\\
\bottomrule
\end{tabular}
\begin{uakeybox}
La pregunta no es «¿qué tecnología es mejor?», sino «¿qué error es inaceptable para esta operación?».
\end{uakeybox}
\end{frame}

% ID: C08-MAP
\begin{frame}{Mapa de la clase}
\UASlideTag{NÚCLEO}{UAIngNavy}
\small
\begin{enumerate}
\item \textbf{Observar}: fallas parciales y conocimiento local.
\item \textbf{Describir}: causalidad y relojes lógicos.
\item \textbf{Decidir}: máquinas de estado, Raft y Paxos.
\item \textbf{Elegir durante una partición}: CAP y PACELC.
\item \textbf{Conservar o evitar conflictos}: Dynamo y Spanner.
\end{enumerate}
\end{frame}

\UASectionSlide{Modelo mínimo y caso conductor}

% ID: C08-M01
\begin{frame}{Modelo mínimo: qué contiene una réplica}
\UASlideTag{NÚCLEO}{UAIngNavy}
\small
\begin{center}
\begin{tikzpicture}[node distance=0.35cm]
\node[uaOrange,text width=2.0cm] (node) {Réplica N1};
\node[uaSoft,text width=1.65cm,right=of node,yshift=1.05cm] (proc) {proceso};
\node[uaSoft,text width=1.65cm,right=of node,yshift=0.35cm] (ram) {RAM};
\node[uaSoft,text width=1.65cm,right=of node,yshift=-0.35cm] (disk) {estado persistente};
\node[uaSoft,text width=1.65cm,right=of node,yshift=-1.05cm] (nic) {red};
\draw[uaArrow] (node.east)--(proc.west);
\draw[uaArrow] (node.east)--(ram.west);
\draw[uaArrow] (node.east)--(disk.west);
\draw[uaArrow] (node.east)--(nic.west);
\end{tikzpicture}
\end{center}
\begin{uakeybox}
Un nodo puede estar sano localmente y, aun así, ignorar qué ocurrió en otra región.
\end{uakeybox}
\end{frame}

% ID: C08-M02
\begin{frame}{Tres situaciones que no debemos confundir}
\UASlideTag{NÚCLEO}{UAIngNavy}
\small
\begin{columns}[T]
\column{0.31\textwidth}
\begin{block}{Sano y comunicado}
Ejecuta y recibe mensajes.
\end{block}
\column{0.31\textwidth}
\begin{block}{Sano pero aislado}
Ejecuta, pero sus mensajes no cruzan a tiempo.
\end{block}
\column{0.31\textwidth}
\begin{block}{Caído}
No ejecuta pasos ni responde.
\end{block}
\end{columns}
\begin{uainfo}
Desde otro nodo, un timeout puede ser compatible con cualquiera de estas historias.
\end{uainfo}
\end{frame}

% ID: C08-M03
\begin{frame}{Timeout: observación local, no diagnóstico}
\UASlideTag{NÚCLEO}{UAIngNavy}
\small
\begin{center}
\begin{tikzpicture}[node distance=0.8cm]
\node[uaOrange,text width=1.4cm] (a) {A};
\node[uaSoft,text width=1.4cm,right=6.8cm of a] (b) {B};
\draw[uaArrow] (a.east)--node[above,uaTiny]{request}(b.west);
\draw[uaArrow,dashed] (b.west)--node[below,uaTiny]{¿respuesta?}(a.east);
\end{tikzpicture}
\end{center}
\begin{itemize}
\item la solicitud nunca llegó;
\item B está lento o pausado;
\item B ejecutó y la respuesta se perdió;
\item la red está particionada.
\end{itemize}
\begin{uakeybox}
El timeout solo prueba que A no obtuvo respuesta dentro de su plazo.
\end{uakeybox}
\end{frame}

% ID: C08-M04
\begin{frame}{Caso conductor: Festival Austral Global}
\UASlideTag{NÚCLEO}{UAIngNavy}
\small
\begin{columns}[T]
\column{0.45\textwidth}
\begin{block}{Recurso}
\texttt{VIP-01}, stock = 1.
\end{block}
\begin{block}{Clientes}
Sofía compra en Buenos Aires; Tomás compra en Madrid.
\end{block}
\column{0.50\textwidth}
\begin{block}{Servicio replicado}
Cinco réplicas: N1, N2, N3, N4 y N5.
\end{block}
\begin{alertblock}{Invariante}
\[
\text{reservas confirmadas para VIP-01}\le 1
\]
\end{alertblock}
\end{columns}
\end{frame}

% ID: C08-M05
\begin{frame}{Predicción inicial}
\UASlideTag{ACTIVIDAD}{UAIngOrange}
\small
La red se divide:
\[
\{N1,N2,N3\}\;|\;\{N4,N5\}
\]
Sofía compra desde el lado de tres nodos y Tomás desde el lado de dos.
\vspace{0.3cm}
\begin{block}{Elija antes de conocer el algoritmo}
¿Qué debería responder el lado minoritario a Tomás?
\begin{enumerate}[A.]
\item Confirmada.
\item Rechazada temporalmente.
\item Pendiente hasta recuperar coordinación.
\end{enumerate}
\end{block}
\end{frame}

\UASectionSlide{Causalidad y relojes lógicos}

% ID: C08-C01
\begin{frame}{Happened-before: tres reglas}
\UASlideTag{NÚCLEO}{UAIngNavy}
\small
\begin{enumerate}
\item Si dos eventos ocurren en el mismo proceso, su orden local induce causalidad.
\item El envío de un mensaje ocurre antes que la recepción de ese mensaje.
\item La relación es transitiva.
\end{enumerate}
\[
a \rightarrow b \wedge b \rightarrow c \Rightarrow a \rightarrow c
\]
significa que existe una cadena de influencia posible desde $a$ hasta $c$.
\end{frame}

% ID: C08-C02

\begin{frame}{Las dos reservas nacen en ramas independientes y llegan al líder}
\UASlideTag{NÚCLEO}{UAIngNavy}
\scriptsize

\begin{center}
\begin{tikzpicture}[x=1.05cm,y=0.86cm]

% ---------------------------------------------------------------------------
% Líneas de proceso
% ---------------------------------------------------------------------------
\draw[thick,draw=UAIngNavy] (0,2)--(8.4,2);
\node[left,font=\scriptsize\bfseries,text=UAIngNavy] at (0,2)
  {BA};

\draw[thick,draw=UAIngNavy] (0,0)--(8.4,0);
\node[left,font=\scriptsize\bfseries,text=UAIngNavy] at (0,0)
  {Líder N1};

\draw[thick,draw=UAIngNavy] (0,-2)--(8.4,-2);
\node[left,font=\scriptsize\bfseries,text=UAIngNavy] at (0,-2)
  {Madrid};

% ---------------------------------------------------------------------------
% Eventos de Buenos Aires
% ---------------------------------------------------------------------------
\node[uaDot] (a1) at (0.9,2) {};
\node[uaTiny,anchor=south] at ([yshift=0.07cm]a1.north)
  {Sofía hace click\\$(a_1)$};

\node[uaDot] (a2) at (2.7,2) {};
\node[uaTiny,anchor=south] at ([yshift=0.07cm]a2.north)
  {envía reserva\\$(a_2)$};

% ---------------------------------------------------------------------------
% Eventos de Madrid
% ---------------------------------------------------------------------------
\node[uaDot] (m1) at (1.3,-2) {};
\node[uaTiny,anchor=north] at ([yshift=-0.07cm]m1.south)
  {Tomás hace click\\$(m_1)$};

\node[uaDot] (m2) at (4.7,-2) {};
\node[uaTiny,anchor=north] at ([yshift=-0.07cm]m2.south)
  {envía reserva\\$(m_2)$};

% ---------------------------------------------------------------------------
% Eventos del líder
% ---------------------------------------------------------------------------
\node[uaDot,fill=UAIngOrange] (l1) at (3.7,0) {};
\node[uaTiny,anchor=north] at ([yshift=-0.07cm]l1.south)
  {recibe a Sofía\\$(l_1)$};

\node[uaDot,fill=UAIngOrange] (l2) at (5.0,0) {};
\node[uaTiny,anchor=north] at ([yshift=-0.07cm]l2.south)
  {append local\\$(l_2)$};

\node[uaDot,fill=UAIngOrange] (l3) at (6.7,0) {};
\node[uaTiny,anchor=north] at ([yshift=-0.07cm]l3.south)
  {recibe a Tomás\\$(l_3)$};

% ---------------------------------------------------------------------------
% Mensajes
% ---------------------------------------------------------------------------
\draw[->,very thick,>=Latex,draw=UAIngOrange]
  (a2) --
  node[uaTiny,fill=white,inner sep=1.2pt,
       pos=0.48,sloped,above]
       {\texttt{Reserve(Sofía)}}
  (l1);

\draw[->,very thick,>=Latex,draw=UAIngOrange]
  (m2) --
  node[uaTiny,fill=white,inner sep=1.2pt,
       pos=0.50,sloped,below]
       {\texttt{Reserve(Tomás)}}
  (l3);

\end{tikzpicture}
\end{center}

\vspace{0.08cm}

\begin{columns}[T]
\column{0.47\textwidth}

\begin{block}{Relaciones causales}
\[
a_1\rightarrow a_2\rightarrow l_1
\]

\[
m_1\rightarrow m_2\rightarrow l_3
\]
\end{block}

\column{0.49\textwidth}

\begin{block}{Concurrencia inicial}
\[
a_2\parallel m_2
\]

Cuando ambas reservas se envían, ninguna contiene información sobre la otra.
\end{block}
\end{columns}

\begin{uakeybox}
El líder ordena sus propios eventos
\(l_1\rightarrow l_2\rightarrow l_3\), pero ese orden de recepción
no convierte a \(a_2\) en causa de \(m_2\), ni a \(m_2\) en causa de \(a_2\).
\end{uakeybox}

\end{frame}

% ID: C08-C03
\begin{frame}{Lamport: convertir relaciones causales en números crecientes}
\UASlideTag{NÚCLEO}{UAIngNavy}
\small
\begin{block}{Variables}
$L$ es el contador lógico local. $t_m$ es la marca lógica incluida por el emisor en el mensaje.
\end{block}
Al recibir:
\[
L\leftarrow \max(L,t_m)+1
\]

\begin{itemize}
\item $\max$ conserva la historia local y la comunicada;
\item $+1$ coloca la recepción estrictamente después del envío.
\end{itemize}
\begin{uainfo}
Ejemplo: $L=4$, $t_m=7$ $\Rightarrow$ nuevo $L=8$.
\end{uainfo}
\end{frame}


% =============================================================================
% Insertar después de la explicación de la regla de Lamport
% =============================================================================

\begin{frame}{La misma ejecución con relojes de Lamport}
\UASlideTag{NÚCLEO}{UAIngNavy}
\scriptsize

\begin{center}
\begin{tikzpicture}[x=1.05cm,y=0.86cm]

% ---------------------------------------------------------------------------
% Líneas de proceso
% ---------------------------------------------------------------------------
\draw[thick,draw=UAIngNavy] (0,2)--(8.4,2);
\node[left,font=\scriptsize\bfseries,text=UAIngNavy] at (0,2)
  {BA};

\draw[thick,draw=UAIngNavy] (0,0)--(8.4,0);
\node[left,font=\scriptsize\bfseries,text=UAIngNavy] at (0,0)
  {Líder N1};

\draw[thick,draw=UAIngNavy] (0,-2)--(8.4,-2);
\node[left,font=\scriptsize\bfseries,text=UAIngNavy] at (0,-2)
  {Madrid};

% ---------------------------------------------------------------------------
% Buenos Aires: L = 1, 2
% ---------------------------------------------------------------------------
\node[uaDot] (a1) at (0.9,2) {};
\node[uaTiny,anchor=south] at ([yshift=0.07cm]a1.north)
  {click \(a_1\)\\\(L=1\)};

\node[uaDot] (a2) at (2.7,2) {};
\node[uaTiny,anchor=south] at ([yshift=0.07cm]a2.north)
  {send \(a_2\)\\\(L=2\)};

% ---------------------------------------------------------------------------
% Madrid: L = 1, 2
% ---------------------------------------------------------------------------
\node[uaDot] (m1) at (1.3,-2) {};
\node[uaTiny,anchor=north] at ([yshift=-0.07cm]m1.south)
  {click \(m_1\)\\\(L=1\)};

\node[uaDot] (m2) at (4.7,-2) {};
\node[uaTiny,anchor=north] at ([yshift=-0.07cm]m2.south)
  {send \(m_2\)\\\(L=2\)};

% ---------------------------------------------------------------------------
% Líder: recibe 2, evento local, recibe 2
% ---------------------------------------------------------------------------
\node[uaDot,fill=UAIngOrange] (l1) at (3.7,0) {};
\node[uaTiny,anchor=north] at ([yshift=-0.07cm]l1.south)
  {receive \(l_1\)\\\(L=3\)};

\node[uaDot,fill=UAIngOrange] (l2) at (5.0,0) {};
\node[uaTiny,anchor=north] at ([yshift=-0.07cm]l2.south)
  {append \(l_2\)\\\(L=4\)};

\node[uaDot,fill=UAIngOrange] (l3) at (6.7,0) {};
\node[uaTiny,anchor=north] at ([yshift=-0.07cm]l3.south)
  {receive \(l_3\)\\\(L=5\)};

% ---------------------------------------------------------------------------
% Mensajes con timestamp t_m = 2
% ---------------------------------------------------------------------------
\draw[->,very thick,>=Latex,draw=UAIngOrange]
  (a2) --
  node[uaTiny,fill=white,inner sep=1.2pt,
       pos=0.48,sloped,above]
       {\(t_m=2\)}
  (l1);

\draw[->,very thick,>=Latex,draw=UAIngOrange]
  (m2) --
  node[uaTiny,fill=white,inner sep=1.2pt,
       pos=0.50,sloped,below]
       {\(t_m=2\)}
  (l3);

\end{tikzpicture}
\end{center}

\vspace{0.08cm}

\begin{uainfo}
El líder comienza con \(L=0\):

\[
\begin{aligned}
L(l_1) &= \max(0,2)+1 = 3,\\
L(l_2) &= 3+1 = 4,\\
L(l_3) &= \max(4,2)+1 = 5.
\end{aligned}
\]
\end{uainfo}

\begin{uakeybox}
Los dos mensajes salen con \(L=2\). La igualdad de timestamps no significa
que sean el mismo evento ni que exista causalidad entre ellos.
Lamport garantiza que \(x\rightarrow y\) implica \(L(x)<L(y)\),
pero el recíproco no es necesariamente cierto.
\end{uakeybox}

\end{frame}


% ID: C08-C04
\begin{frame}{Ejemplo Lamport sobre las dos reservas}
\UASlideTag{ACTIVIDAD}{UAIngOrange}
\small
\begin{tabular}{p{0.33\textwidth}p{0.25\textwidth}p{0.25\textwidth}}
\toprule
\textbf{Evento} & \textbf{Proceso} & \textbf{Lamport}\\
\midrule
Sofía presiona comprar & BA & 1\\
BA envía reserva & BA & 2\\
Tomás presiona comprar & Madrid & 1\\
Madrid envía reserva & Madrid & 2\\
Líder recibe a Sofía & líder & 3\\
Líder hace append & líder & 4\\
Líder recibe a Tomás & líder & 5\\
\bottomrule
\end{tabular}
\end{frame}

% ID: C08-C05
\begin{frame}{Límite de Lamport}
\UASlideTag{NÚCLEO}{UAIngNavy}
\small
\[
a\rightarrow b \Longrightarrow L(a)<L(b)
\]
pero:
\[
L(a)<L(b) \not\Longrightarrow a\rightarrow b.
\]
\begin{block}{Por qué}
Los números de Lamport permiten respetar causalidad y construir un desempate total convencional, pero no distinguen por sí solos una dependencia real de una simple concurrencia.
\end{block}
\end{frame}

% ID: C08-C06
\begin{frame}{Relojes vectoriales: registrar conocimiento por actor}
\UASlideTag{NÚCLEO}{UAIngNavy}
\small
Usamos el orden:
\[
[\text{BA},\ \text{Madrid},\ \text{Líder}]
\]
Cada proceso incrementa su componente y, al recibir, toma el máximo componente a componente antes de incrementar el suyo.
\begin{block}{Comparación}
$V_a<V_b$ si todos los componentes de $V_a$ son menores o iguales y al menos uno es estrictamente menor. Si ninguno domina al otro, son concurrentes.
\end{block}
\end{frame}

% =============================================================================
% Insertar después de las reglas de relojes vectoriales
% Orden de componentes: [BA, Madrid, Líder]
% =============================================================================

% =============================================================================
% La misma ejecución con relojes vectoriales
% =============================================================================

\begin{frame}{La misma ejecución con relojes vectoriales}
\UASlideTag{NÚCLEO}{UAIngNavy}
\scriptsize

\begin{center}
\begin{tikzpicture}[x=1.02cm,y=0.72cm]

% ---------------------------------------------------------------------------
% Líneas de proceso
% ---------------------------------------------------------------------------
\draw[thick,draw=UAIngNavy] (0,1.35)--(8.1,1.35);
\node[left,font=\scriptsize\bfseries,text=UAIngNavy]
  at (0,1.65) {BA};

\draw[thick,draw=UAIngNavy] (0,0)--(8.1,0);
\node[left,font=\scriptsize\bfseries,text=UAIngNavy]
  at (0,0) {Líder N1};

\draw[thick,draw=UAIngNavy] (0,-1.35)--(8.1,-1.35);
\node[left,font=\scriptsize\bfseries,text=UAIngNavy]
  at (0,-1.65) {Madrid};

% ---------------------------------------------------------------------------
% Eventos de Buenos Aires
% ---------------------------------------------------------------------------
\node[circle,fill=UAIngNavy,inner sep=2.2pt] (a1) at (0.9,1.35) {};
\node[font=\tiny,align=center,anchor=south]
  at ([yshift=0.06cm]a1.north)
  {\(a_1\): click\\\([1,0,0]\)};

\node[circle,fill=UAIngNavy,inner sep=2.2pt] (a2) at (3.0,1.35) {};
\node[font=\tiny,align=center,anchor=south]
  at ([yshift=0.06cm]a2.north)
  {\(a_2\): send\\\([2,0,0]\)};

% ---------------------------------------------------------------------------
% Eventos de Madrid
% ---------------------------------------------------------------------------
\node[circle,fill=UAIngNavy,inner sep=2.2pt] (m1) at (1.3,-1.35) {};
\node[font=\tiny,align=center,anchor=north]
  at ([yshift=-0.06cm]m1.south)
  {\(m_1\): click\\\([0,1,0]\)};

\node[circle,fill=UAIngNavy,inner sep=2.2pt] (m2) at (6.5,-1.35) {};
\node[font=\tiny,align=center,anchor=north]
  at ([yshift=-0.06cm]m2.south)
  {\(m_2\): send\\\([0,2,0]\)};

% ---------------------------------------------------------------------------
% Eventos del líder
% ---------------------------------------------------------------------------
\node[circle,fill=UAIngOrange,inner sep=2.4pt] (l1) at (3.0,0) {};
\node[font=\tiny,align=center,anchor=north]
  at ([yshift=-0.08cm]l1.south)
  {\(l_1\): recibe a Sofía\\\([2,0,1]\)};

\node[circle,fill=UAIngOrange,inner sep=2.4pt] (l2) at (4.7,0) {};
\node[font=\tiny,align=center,anchor=north]
  at ([yshift=-0.08cm]l2.south)
  {\(l_2\): append\\\([2,0,2]\)};

\node[circle,fill=UAIngOrange,inner sep=2.4pt] (l3) at (6.5,0) {};
\node[font=\tiny,align=center,anchor=south]
  at ([yshift=0.08cm]l3.north)
  {\(l_3\): recibe a Tomás\\\([2,2,3]\)};

% ---------------------------------------------------------------------------
% Mensajes
% ---------------------------------------------------------------------------
\draw[->,very thick,>=Latex,draw=UAIngOrange]
  (a2)--(l1);

\draw[->,very thick,>=Latex,draw=UAIngOrange]
  (m2)--(l3);

\end{tikzpicture}
\end{center}

\vspace{-0.12cm}

\begin{columns}[T,onlytextwidth]

\column{0.47\textwidth}

\begin{uainfo}
\textbf{Orden de componentes}

\[
[\text{BA},\text{Madrid},\text{Líder}]
\]

Cada proceso incrementa solamente su propio componente.
\end{uainfo}

\column{0.49\textwidth}

\begin{uakeybox}
\[
[2,0,0]\parallel[0,2,0]
\]

Los vectores de las dos reservas son incomparables: ninguna reserva contiene
conocimiento causal de la otra.
\end{uakeybox}

\end{columns}

\end{frame}

% =============================================================================
% Cálculo paso a paso del vector del líder
% =============================================================================

\begin{frame}{Cómo actualiza su reloj vectorial el líder}
\UASlideTag{NÚCLEO}{UAIngNavy}
\footnotesize

\begin{block}{Regla de recepción}
Al recibir un mensaje, el proceso:

\begin{enumerate}
\item calcula el máximo componente a componente entre su vector y el recibido;
\item incrementa su propio componente para representar el nuevo evento local.
\end{enumerate}
\end{block}

\vspace{-0.25cm}

\[
\begin{aligned}
V(l_1)
  &= \max\bigl([0,0,0],[2,0,0]\bigr)+[0,0,1] = [2,0,1],
\\
V(l_2)
  &= [2,0,1]+[0,0,1]= [2,0,2],
\\
V(l_3)
  &= \max\bigl([2,0,2],[0,2,0]\bigr)+[0,0,1]= [2,2,3].
\end{aligned}
\]

\vspace{-0.15cm}

\begin{columns}[T,onlytextwidth]

\column{0.47\textwidth}

\begin{uainfo}
\footnotesize
\textbf{¿Qué expresa \([2,2,3]\)? El líder conoce:}

\begin{itemize}
\item dos eventos de Buenos Aires;
\item dos eventos de Madrid;
\item tres eventos locales propios.
\end{itemize}
\end{uainfo}

\column{0.49\textwidth}
\begin{uakeybox}
\footnotesize
El resultado es \([2,2,3]\), no \([2,2,5]\).

El valor \(5\) pertenece al reloj escalar de Lamport. En el reloj vectorial,
el tercer componente cuenta solamente los eventos locales del líder.
\end{uakeybox}

\end{columns}

\end{frame}

% ID: C08-C07
\begin{frame}{Vector clocks en la historia}
\UASlideTag{NÚCLEO}{UAIngNavy}
\small
\begin{tabular}{p{0.43\textwidth}p{0.35\textwidth}}
\toprule
\textbf{Evento} & \textbf{Vector}\\
\midrule
BA envía reserva de Sofía & $[2,0,0]$\\
Madrid envía reserva de Tomás & $[0,2,0]$\\
Líder recibe a Sofía & $[2,0,1]$\\
Líder hace append & $[2,0,2]$\\
Líder recibe a Tomás & $[2,2,3]$\\
\bottomrule
\end{tabular}
\begin{uawarn}
El reloj escalar del líder termina en $L=5$, pero el tercer componente vectorial es 3 porque cuenta sus tres eventos locales. No se mezclan ambos relojes.
\end{uawarn}
\end{frame}

% ID: C08-C08
\begin{frame}{Tres relojes, tres preguntas}
\UASlideTag{NÚCLEO}{UAIngNavy}
\scriptsize
\begin{tabular}{p{0.22\textwidth}p{0.30\textwidth}p{0.34\textwidth}}
\toprule
\textbf{Reloj} & \textbf{Pregunta} & \textbf{Uso}\\
\midrule
pared & ¿qué fecha y hora civil es? & auditoría, timestamps humanos\\
monotónico & ¿cuánto duró una espera? & timeouts, latencias, leases\\
lógico & ¿qué orden causal debe respetarse? & causalidad, versiones, logs\\
\bottomrule
\end{tabular}
\begin{uakeybox}
Los relojes describen relaciones. Todavía no deciden quién obtiene VIP-01.
\end{uakeybox}
\end{frame}

\UASectionSlide{Máquinas de estado y consenso}

% ID: C08-S01
\begin{frame}{Máquina de estados: una regla determinista}
\UASlideTag{NÚCLEO}{UAIngNavy}
\small
\[
(estado,comando)\longrightarrow(nuevo\ estado,respuesta)
\]
\begin{center}
\begin{tikzpicture}[node distance=0.55cm]
\node[uaOrange,text width=2.0cm] (avail) {AVAILABLE};
\node[uaSoft,text width=2.25cm,right=2.2cm of avail,yshift=0.75cm] (sof) {RESERVED\\BY SOFÍA};
\node[uaSoft,text width=2.25cm,right=2.2cm of avail,yshift=-0.75cm] (tom) {RESERVED\\BY TOMÁS};
\draw[uaOrangeArrow] (avail.east)--node[above,uaTiny]{Reserve(Sofía)}(sof.west);
\draw[uaArrow] (avail.east)--node[below,uaTiny]{Reserve(Tomás)}(tom.west);
\end{tikzpicture}
\end{center}
Una reserva posterior sobre un estado ya reservado se rechaza sin cambiar el estado.
\end{frame}

% ID: C08-S02
\begin{frame}{Máquina de estados replicada}
\UASlideTag{NÚCLEO}{UAIngNavy}
\small
Cada réplica contiene:
\begin{itemize}
\item una copia del log de comandos;
\item el mismo código determinista;
\item una copia local del estado.
\end{itemize}
\begin{uakeybox}
Si todas aplican los mismos comandos comprometidos en el mismo orden, convergen al mismo estado.
\end{uakeybox}
\end{frame}

% ID: C08-S03
\begin{frame}{Safety y liveness}
\UASlideTag{NÚCLEO}{UAIngNavy}
\small
\begin{columns}[T]
\column{0.48\textwidth}
\begin{block}{Safety: nada incorrecto ocurre}
\begin{itemize}
\item no hay dos comandos distintos comprometidos en el mismo índice;
\item una entrada committed no se reemplaza;
\item VIP-01 no se confirma para dos personas.
\end{itemize}
\end{block}
\column{0.48\textwidth}
\begin{block}{Liveness: el sistema vuelve a progresar}
Requiere mayoría comunicada, líder estable, procesos activos y entrega eventual de los mensajes necesarios.
\end{block}
\end{columns}
\end{frame}

% ID: C08-S04
\begin{frame}{Raft: término, índice y entrada de log}
\UASlideTag{NÚCLEO}{UAIngNavy}
\small
\begin{itemize}
\item \textbf{Término}: número creciente que identifica una ronda de elección y liderazgo.
\item \textbf{Índice}: posición estable de una entrada dentro del log.
\item \textbf{Entrada}: par $(term,command)$ en una posición.
\end{itemize}
\begin{block}{Ejemplo}
Los índices 1--86 ya existen. El líder del término 12 agrega:
\[
index=87,\quad term=12,\quad command=\operatorname{Reserve}(\text{VIP-01},\text{Sofía})
\]
\end{block}
\end{frame}

% ID: C08-S05
\begin{frame}{Mensaje o líder obsoleto}
\UASlideTag{NÚCLEO}{UAIngNavy}
\small
Un mensaje es obsoleto cuando contiene un término menor que el término ya conocido por el receptor.
\begin{center}
\begin{tikzpicture}[node distance=1.1cm]
\node[uaSoft,text width=2.1cm] (n1) {N1 despierta\\term 12};
\node[uaOrange,text width=2.1cm,right=of n1] (n2) {N2 líder\\term 13};
\draw[uaArrow] (n1)--node[above,uaTiny]{AppendEntries term 12}(n2);
\node[below=0.5cm of n2,uaTiny,text=UAIngOrange] {rechazado};
\end{tikzpicture}
\end{center}
\end{frame}

% ID: C08-S06
\begin{frame}{Raft paso 1: la propuesta existe solo en N1}
\UASlideTag{NÚCLEO}{UAIngNavy}
\small
\begin{tabular}{p{0.13\textwidth}p{0.58\textwidth}p{0.18\textwidth}}
\toprule
\textbf{Nodo} & \textbf{Última entrada} & \textbf{Estado}\\
\midrule
N1 & 87 / term 12 / Reserve(Sofía) & local\\
N2 & hasta 86 & ---\\
N3 & hasta 86 & ---\\
N4 & hasta 86 & ---\\
N5 & hasta 86 & ---\\
\bottomrule
\end{tabular}
\begin{alertblock}{Todavía no}
Una propuesta local no autoriza responder «compra confirmada».
\end{alertblock}
\end{frame}

% ID: C08-S07
\begin{frame}{Raft paso 2: dos copias todavía no son mayoría}
\UASlideTag{NÚCLEO}{UAIngNavy}
\small
\begin{tabular}{p{0.13\textwidth}p{0.58\textwidth}p{0.18\textwidth}}
\toprule
\textbf{Nodo} & \textbf{Última entrada} & \textbf{Estado}\\
\midrule
N1 & 87 / term 12 / Reserve(Sofía) & local\\
N2 & 87 / term 12 / Reserve(Sofía) & local\\
N3 & hasta 86 & ---\\
N4 & hasta 86 & ---\\
N5 & hasta 86 & ---\\
\bottomrule
\end{tabular}
N2 persiste la entrada y responde al líder.
\[
2/5<3/5
\]
\begin{uakeybox}
La seguridad no proviene solo del número de copias, sino de una mayoría intersectante y de las reglas del protocolo.
\end{uakeybox}
\end{frame}

% ID: C08-S08
\begin{frame}{Raft paso 3: la tercera réplica cambia la situación}
\UASlideTag{NÚCLEO}{UAIngNavy}
\small
\begin{tabular}{p{0.13\textwidth}p{0.58\textwidth}p{0.18\textwidth}}
\toprule
\textbf{Nodo} & \textbf{Última entrada} & \textbf{Estado}\\
\midrule
N1 & 87 / term 12 / Reserve(Sofía) & local\\
N2 & 87 / term 12 / Reserve(Sofía) & local\\
N3 & 87 / term 12 / Reserve(Sofía) & local\\
N4 & hasta 86 & ---\\
N5 & hasta 86 & ---\\
\bottomrule
\end{tabular}
N3 persiste la misma entrada.
\[
3/5
\]
Toda mayoría futura de tres nodos debe compartir al menos un nodo con $\{N1,N2,N3\}$.
\begin{uakeybox}
La entrada puede avanzar a committed porque alcanzó la mayoría exigida por el protocolo.
\end{uakeybox}
\end{frame}

% ID: C08-S09
\begin{frame}{Raft paso 4: quién compromete, aplica y responde}
\UASlideTag{NÚCLEO}{UAIngNavy}
\small
\begin{enumerate}
\item N2 y N3 confirman al líder que persistieron la entrada.
\item N1 detecta mayoría y actualiza \texttt{commitIndex=87}.
\item N1 aplica la entrada en \textbf{su propia} máquina de estados.
\item N1 responde «compra confirmada» a Sofía.
\item N2 y N3 reciben \texttt{leaderCommit=87}, avanzan y aplican localmente.
\end{enumerate}
\begin{uawarn}
Replicada $\neq$ committed $\neq$ aplicada $\neq$ confirmada al cliente.
\end{uawarn}
\end{frame}

% ID: C08-S10
\begin{frame}{Crash antes y después de la mayoría}
\UASlideTag{ACTIVIDAD}{UAIngOrange}
\small
\begin{columns}[T]
\column{0.48\textwidth}
\begin{block}{Antes de mayoría}
N1 contiene la entrada y cae. No estaba committed y puede desaparecer. Sofía no debió recibir éxito.
\end{block}
\column{0.48\textwidth}
\begin{block}{Después de mayoría}
N1, N2 y N3 persistieron; N1 confirmó y cae. N2 o N3 conserva la decisión comprometida.
\end{block}
\end{columns}
\end{frame}

% ID: C08-S11
\begin{frame}{Total-order broadcast: por qué aparece aquí}
\UASlideTag{NÚCLEO}{UAIngNavy}
\small
\begin{block}{Necesidad}
Máquinas deterministas solo convergen si reciben los mismos comandos comprometidos y los aplican en el mismo orden.
\end{block}
\begin{block}{Relación con consenso}
Consenso elige un valor para una posición. Repetir consenso para 87, 88, 89, ... construye un log ordenado; ese servicio implementa total-order broadcast.
\end{block}
\begin{uainfo}
Un timestamp Lamport etiqueta eventos, pero no garantiza que todas las réplicas reciban el mismo conjunto de comandos ni que no existan huecos.
\end{uainfo}
\end{frame}

% ID: C08-S12
\begin{frame}{Pausa de proceso, GC y fencing}
\UASlideTag{NÚCLEO}{UAIngNavy}
\small
\begin{itemize}
\item \textbf{GC}: pausa del runtime para recuperar memoria; el proceso puede dejar de ejecutar sin estar muerto.
\item \textbf{Fencing token}: número creciente que un recurso externo usa para rechazar operaciones de una autoridad antigua.
\end{itemize}
\begin{center}
N1 líder term 12 $\rightarrow$ pausa $\rightarrow$ N2 líder term 13 $\rightarrow$ N1 despierta obsoleto
\end{center}
\begin{uakeybox}
El término protege al cluster; un recurso externo puede requerir además un token de fencing.
\end{uakeybox}
\end{frame}

% ID: C08-S13
\begin{frame}{Paxos: misma meta, organización diferente}
\UASlideTag{EXTENSIÓN}{UAIngOrange}
\scriptsize
\begin{tabular}{p{0.20\textwidth}p{0.31\textwidth}p{0.31\textwidth}}
\toprule
\textbf{Paso} & \textbf{Raft} & \textbf{Multi-Paxos}\\
\midrule
liderazgo & líder, term 12 & proposer estable, ballot 12\\
posición & index 87 & slot 87\\
propuesta & AppendEntries & Accept\\
réplicas & followers & acceptors\\
decisión & mayoría replica y commit & mayoría acepta, valor elegido\\
aplicación & máquina de estados & learner y máquina de estados\\
\bottomrule
\end{tabular}
\begin{uakeybox}
Ambos impiden decisiones incompatibles para la misma posición, aunque organizan roles y mensajes de forma diferente.
\end{uakeybox}
\end{frame}

% ID: C08-S14
\begin{frame}{Caso real: dónde aparecen estas ideas en Kubernetes}
\UASlideTag{CASO REAL}{UAIngOrange}
\small
\begin{center}
\texttt{kubectl} $\rightarrow$ API server $\rightarrow$ líder etcd $\rightarrow$ seguidores etcd
\end{center}
\begin{tabular}{p{0.30\textwidth}p{0.57\textwidth}}
\toprule
\textbf{Concepto} & \textbf{Aplicación en etcd}\\
\midrule
líder y término & un líder Raft atiende escrituras y rechaza términos viejos\\
log e índice & cada cambio ocupa una posición del log replicado\\
mayoría y commit & se confirma tras alcanzar quorum\\
máquina de estados & todos obtienen el mismo key--value store\\
WAL/snapshot & cada miembro recuerda decisiones tras reiniciar\\
\bottomrule
\end{tabular}
\end{frame}

% ID: C08-S15
\begin{frame}{Tres logs que no debemos confundir}
\UASlideTag{NÚCLEO}{UAIngNavy}
\scriptsize
\begin{tabular}{p{0.24\textwidth}p{0.30\textwidth}p{0.34\textwidth}}
\toprule
\textbf{Log} & \textbf{Pregunta} & \textbf{Ejemplo}\\
\midrule
consenso & ¿qué comando se acordó y en qué orden? & log Raft\\
WAL local & ¿qué recuerda físicamente el nodo tras un crash? & WAL de la DB/etcd\\
eventos & ¿qué hechos pueden leer y reprocesar otros componentes? & Kafka\\
\bottomrule
\end{tabular}
\end{frame}

\UASectionSlide{CAP, PACELC y garantías}

% ID: C08-P01
\begin{frame}{La partición ocurre después del commit de Sofía}
\UASlideTag{NÚCLEO}{UAIngNavy}
\small
\[
\{N1,N2,N3\}\;|\;\{N4,N5\}
\]
La mayoría conoce la reserva committed de Sofía. La minoría está sana localmente, pero no recibe ese hecho. Tomás intenta comprar desde Madrid.
\end{frame}

% ID: C08-P02
\begin{frame}{Dos historias indistinguibles para la minoría}
\UASlideTag{NÚCLEO}{UAIngNavy}
\small
\begin{columns}[T]
\column{0.48\textwidth}
\begin{block}{H1}
Sofía ya fue confirmada, pero el mensaje no cruza la partición.
\end{block}
\column{0.48\textwidth}
\begin{block}{H2}
Sofía nunca fue confirmada.
\end{block}
\end{columns}
N4 y N5 observan lo mismo en ambas historias. Una respuesta local idéntica puede ser correcta en una e incorrecta en la otra.
\end{frame}

% ID: C08-P03
\begin{frame}{CAP: la decisión durante una partición}
\UASlideTag{NÚCLEO}{UAIngNavy}
\small
\begin{itemize}
\item \textbf{Consistencia fuerte}: una sola historia compatible con una copia autoritativa.
\item \textbf{Disponibilidad CAP}: cada solicitud a un nodo no fallado termina con una respuesta válida de la operación.
\item \textbf{Partición}: grupos de nodos vivos no pueden intercambiar los mensajes necesarios a tiempo.
\end{itemize}
\begin{uakeybox}
Durante la partición, confirmar localmente a Tomás arriesga doble venta; esperar o rechazar preserva el invariante, pero sacrifica disponibilidad CAP para esa operación.
\end{uakeybox}
\end{frame}

% ID: C08-P04
\begin{frame}{Uptime, disponibilidad CAP y disponibilidad semántica}
\UASlideTag{NÚCLEO}{UAIngNavy}
\scriptsize
\begin{tabular}{p{0.25\textwidth}p{0.62\textwidth}}
\toprule
\textbf{Concepto} & \textbf{Qué mide}\\
\midrule
uptime & fracción de una ventana mensual, trimestral o anual en que el proceso estuvo alcanzable\\
disponibilidad CAP & si la operación termina con respuesta válida durante una partición\\
disponibilidad semántica & fracción de intentos válidos que cumplen la promesa de negocio\\
\bottomrule
\end{tabular}
\begin{uainfo}
Un 503 rápido demuestra que el proceso está vivo, pero no satisface disponibilidad CAP para la compra.
\end{uainfo}
\end{frame}

% ID: C08-P05
\begin{frame}{PACELC: el costo continúa cuando la red funciona}
\UASlideTag{NÚCLEO}{UAIngNavy}
\small
\begin{columns}[T]
\column{0.48\textwidth}
\begin{block}{Lectura local}
15 ms, pero puede observar una réplica atrasada.
\end{block}
\column{0.48\textwidth}
\begin{block}{Lectura coordinada}
180 ms, con garantía más fuerte.
\end{block}
\end{columns}
\begin{uakeybox}
PACELC organiza dos preguntas: durante una partición, consistencia o disponibilidad; en régimen normal, latencia o consistencia.
\end{uakeybox}
\end{frame}

% ID: C08-P06
\begin{frame}{Recencia, linearizabilidad y serializabilidad}
\UASlideTag{NÚCLEO}{UAIngNavy}
\scriptsize
\begin{tabular}{p{0.22\textwidth}p{0.68\textwidth}}
\toprule
\textbf{Concepto} & \textbf{Definición breve}\\
\midrule
recencia & qué tan nuevo es el estado observado\\
linearizabilidad & cada operación parece atómica entre invocación y respuesta y respeta el orden real de operaciones no superpuestas\\
serializabilidad & el resultado de transacciones equivale a algún orden serial, sin exigir orden real externo\\
consistencia externa & serializabilidad más respeto del orden observado por los clientes\\
\bottomrule
\end{tabular}
\end{frame}

\UASectionSlide{Quórums, versiones y Dynamo}

% ID: C08-D01
\begin{frame}{Mayoría en consenso y quorum Dynamo-like no son sinónimos}
\UASlideTag{NÚCLEO}{UAIngNavy}
\scriptsize
\begin{tabular}{p{0.22\textwidth}p{0.34\textwidth}p{0.34\textwidth}}
\toprule
 & \textbf{Consenso} & \textbf{Dynamo-like}\\
\midrule
pregunta & qué entrada autoritativa se elige & cuántas réplicas participan\\
unidad & posición/slot del log & lectura o escritura de una key\\
concurrencia & protocolo impide commits incompatibles & puede conservar versiones concurrentes\\
resultado & orden autoritativo & conjunto de versiones y política de merge\\
intersección & parte de reglas de seguridad & propiedad combinatoria de $R,W,N$\\
\bottomrule
\end{tabular}
\end{frame}

% ID: C08-D02
\begin{frame}{Por qué necesitamos intersección}
\UASlideTag{NÚCLEO}{UAIngNavy}
\small
Con $N=5$, una escritura llega a $\{N1,N2\}$ y una lectura posterior consulta $\{N4,N5\}$. Los conjuntos son disjuntos: la lectura puede no encontrar ninguna copia actualizada.
\begin{uakeybox}
Queremos que todo conjunto ideal de lectura comparta al menos un nodo con todo conjunto ideal de escritura.
\end{uakeybox}
\end{frame}

% ID: C08-D03
\begin{frame}{N, R y W: construir la intersección}
\UASlideTag{ACTIVIDAD}{UAIngOrange}
\small
\begin{itemize}
\item $N$: réplicas preferidas para la key.
\item $W$: respuestas requeridas para aceptar la escritura.
\item $R$: respuestas requeridas para completar la lectura.
\end{itemize}
\[
R+W>N
\]
Con $N=5$, $W=3$, $R=3$, los conjuntos ideales deben intersectarse.
\begin{uawarn}
La intersección no crea automáticamente linearizabilidad, líder, orden total ni resolución semántica.
\end{uawarn}
\end{frame}

% ID: C08-D04
\begin{frame}{Contraejemplo 1: escrituras concurrentes}
\UASlideTag{NÚCLEO}{UAIngNavy}
\small
Dos clientes escriben versiones distintas sin conocer la otra. El nodo de intersección puede conservar ambas ramas.
\begin{block}{Por qué contradice la conclusión ingenua}
$R+W>N$ permite descubrir que existen versiones distintas, pero no decide cuál es posterior si fueron concurrentes.
\end{block}
\end{frame}

% ID: C08-D05
\begin{frame}{Contraejemplo 2: sloppy quorum}
\UASlideTag{NÚCLEO}{UAIngNavy}
\small
Durante una falla, la escritura se almacena temporalmente en nodos fuera del conjunto preferido. Antes del hinted handoff, una lectura sobre las réplicas preferidas puede no cruzarse con el conjunto efectivo de escritura.
\begin{block}{Lección}
La fórmula usa conjuntos ideales; la implementación real debe considerar dónde quedaron efectivamente las copias.
\end{block}
\end{frame}

% ID: C08-D06
\begin{frame}{Contraejemplo 3: last-write-wins y relojes desalineados}
\UASlideTag{NÚCLEO}{UAIngNavy}
\small
Una actualización causalmente posterior recibe un timestamp físico menor por desalineación del reloj y puede ser descartada como «más vieja».
\begin{block}{Lección}
Intersección y timestamps físicos no sustituyen un modelo correcto de causalidad.
\end{block}
\end{frame}

% ID: C08-D07
\begin{frame}{Vector clocks reaparecen como genealogía del carrito}
\UASlideTag{NÚCLEO}{UAIngNavy}
\small
Partimos de:
\[
v_0=[0,0]
\]
Durante la partición:
\[
\text{BA agrega libro}\rightarrow[1,0]
\]
\[
\text{Madrid agrega auriculares}\rightarrow[0,1]
\]
Ninguna versión domina a la otra: son ramas concurrentes.
\end{frame}

% ID: C08-D08
\begin{frame}{Carrito fusionable; entrada VIP no fusionable}
\UASlideTag{NÚCLEO}{UAIngNavy}
\small
\begin{columns}[T]
\column{0.48\textwidth}
\begin{block}{Carrito}
Puede conservar libro y auriculares y producir una versión descendiente que incluya ambas intenciones.
\end{block}
\column{0.48\textwidth}
\begin{alertblock}{VIP-01}
No existe un merge válido «Sofía + Tomás». El invariante exige una única decisión autoritativa.
\end{alertblock}
\end{columns}
\end{frame}

% ID: C08-D09
\begin{frame}{Dynamo: localizar, escribir, conservar y reparar}
\UASlideTag{NÚCLEO}{UAIngNavy}
\small
\begin{enumerate}
\item aplicar hash a la key y construir una preference list;
\item escribir en réplicas hasta reunir $W$ respuestas;
\item conservar versiones concurrentes en vez de perderlas;
\item responder según la política de disponibilidad;
\item reparar la divergencia en segundo plano o durante lecturas.
\end{enumerate}
\end{frame}

% ID: C08-D10
\begin{frame}{Sloppy quorum, hint y hinted handoff}
\UASlideTag{NÚCLEO}{UAIngNavy}
\small
\begin{itemize}
\item \textbf{Sloppy quorum}: aceptar temporalmente la copia en un nodo alternativo.
\item \textbf{Hint}: metadato que recuerda cuál era el propietario preferido.
\item \textbf{Hinted handoff}: devolver la copia al propietario cuando se recupera.
\end{itemize}
\begin{uakeybox}
La disponibilidad desplaza costo hacia almacenamiento temporal, red, I/O y reparación posterior.
\end{uakeybox}
\end{frame}

% ID: C08-D11
\begin{frame}{Read repair, anti-entropy y tasa de divergencia}
\UASlideTag{NÚCLEO}{UAIngNavy}
\small
\begin{itemize}
\item \textbf{Read repair}: corrige las réplicas consultadas al detectar versiones atrasadas.
\item \textbf{Anti-entropy}: compara rangos en background para reparar datos fríos.
\item \textbf{Tasa de divergencia}: diferencias nuevas creadas por unidad de tiempo.
\item \textbf{Capacidad de reparación}: diferencias corregidas por unidad de tiempo.
\end{itemize}
\[
D_{nuevo}>R_{resuelto}\Rightarrow\text{la deuda crece}
\]
\end{frame}

% ID: C08-D12
\begin{frame}{Por qué Dynamo aceptó esa complejidad}
\UASlideTag{CASO REAL}{UAIngOrange}
\small
El Dynamo original incorporó consistent hashing, preference lists, quórums, versionado, sloppy quorum, hinted handoff, read repair y anti-entropy para sostener la prioridad \emph{always-on} de algunos servicios de Amazon.
\begin{uakeybox}
No son sinónimos de «complejidad»; cada mecanismo paga una parte concreta de la disponibilidad durante fallas.
\end{uakeybox}
\end{frame}

\UASectionSlide{Spanner, TrueTime y orden externo}

% ID: C08-G01
\begin{frame}{Por qué aparece Spanner en la misma historia}
\UASlideTag{NÚCLEO}{UAIngNavy}
\small
La compra de VIP-01 puede involucrar:
\begin{itemize}
\item grupo de entradas;
\item grupo de pagos;
\item grupo de derechos del usuario.
\end{itemize}
La operación debe ser atómica entre grupos y, si Madrid inicia una acción después de observar la confirmación de Buenos Aires, el orden externo debe respetarse.
\end{frame}

% ID: C08-G02
\begin{frame}{Spanner: consenso localizado y transacción multi-grupo}
\UASlideTag{NÚCLEO}{UAIngNavy}
\small
Cada grupo de datos mantiene replicación y consenso propios. Una transacción que toca varios grupos coordina solo los grupos necesarios y asigna un timestamp de commit común.
\begin{uakeybox}
El objetivo no es un único líder global para toda la base, sino coordinación localizada con una garantía transaccional global.
\end{uakeybox}
\end{frame}

% ID: C08-G03
\begin{frame}{TrueTime: intervalo de tiempo con incertidumbre explícita}
\UASlideTag{NÚCLEO}{UAIngNavy}
\small
Spanner sincroniza servidores con referencias como GPS y relojes atómicos mediante time masters. Entre sincronizaciones existe deriva.
\[
TT.now()=[earliest,latest]
\]
Conceptualmente:
\[
earliest=\widehat{t}-\epsilon,\qquad latest=\widehat{t}+\epsilon
\]
\begin{uainfo}
TrueTime no afirma conocer un instante exacto: expone una cota que incluye error de sincronización y deriva.
\end{uainfo}
\end{frame}

% ID: C08-G04
\begin{frame}{Quién asigna el timestamp de commit}
\UASlideTag{NÚCLEO}{UAIngNavy}
\small
El coordinador de la transacción elige $s$ respetando:
\begin{itemize}
\item timestamps de preparación de los grupos;
\item la cota superior de TrueTime;
\item dependencias y restricciones de la transacción.
\end{itemize}
Ejemplo:
\[
TT.now()=[4,16]\text{ ms},\qquad s=16\text{ ms}
\]
\end{frame}

% ID: C08-G05
\begin{frame}{Commit wait: esperar hasta demostrar que $s$ quedó en el pasado}
\UASlideTag{NÚCLEO}{UAIngNavy}
\small
El sistema espera hasta observar:
\[
TT.now().earliest>s
\]
Por ejemplo:
\[
TT.now()=[17,29]\text{ ms}\Rightarrow 17>16
\]
\begin{uakeybox}
La espera convierte incertidumbre temporal acotada en una garantía de orden externo.
\end{uakeybox}
\end{frame}

% ID: C08-G06
\begin{frame}{Por qué se llama consistencia externa}
\UASlideTag{NÚCLEO}{UAIngNavy}
\small
Si un cliente recibe el commit de $T_1$ y después inicia $T_2$, toda serialización válida debe respetar:
\[
T_1<T_2
\]
Se llama externa porque respeta el orden observado por clientes fuera del motor, no solo algún orden serial internamente válido.
\end{frame}

\UASectionSlide{Integración y cierre}

% ID: C08-I01
\begin{frame}{Cada mecanismo responde una pregunta distinta}
\UASlideTag{NÚCLEO}{UAIngNavy}
\scriptsize
\begin{tabular}{p{0.23\textwidth}p{0.63\textwidth}}
\toprule
\textbf{Mecanismo} & \textbf{Pregunta}\\
\midrule
Lamport/vector & ¿qué eventos están relacionados o son concurrentes?\\
Raft/Paxos & ¿qué comando se convierte en decisión autoritativa?\\
WAL local & ¿qué recuerda cada nodo después de un crash?\\
Kafka & ¿qué hechos pueden leer y reprocesar otros componentes?\\
CAP/PACELC & ¿qué error y qué costo aceptamos bajo red imperfecta?\\
Dynamo & ¿cómo seguir respondiendo y reparar divergencia?\\
Spanner & ¿cómo sostener transacciones globales y orden externo?\\
\bottomrule
\end{tabular}
\end{frame}

% ID: C08-I02
\begin{frame}{Actividad integradora: cuatro operaciones, cuatro decisiones}
\UASlideTag{ACTIVIDAD}{UAIngOrange}
\small
Para cada operación declare invariante, mecanismo, comportamiento bajo partición, costo normal y prueba de falla:
\begin{enumerate}
\item última entrada VIP;
\item carrito compartido;
\item líder del control plane;
\item ranking global que entrega premios.
\end{enumerate}
\end{frame}

% ID: C08-I03
\begin{frame}{Exit ticket}
\UASlideTag{ACTIVIDAD}{UAIngOrange}
\small
Responda en una frase cada punto:
\begin{enumerate}
\item ¿Por qué $L(a)<L(b)$ no prueba causalidad?
\item ¿Cómo detecta un vector que dos eventos son concurrentes?
\item ¿Cuándo queda committed la entrada 87 en Raft?
\item ¿Qué diferencia hay entre log de consenso, WAL y Kafka?
\item ¿Por qué $R+W>N$ no implica linearizabilidad?
\item ¿Qué costo paga Spanner para obtener orden externo?
\end{enumerate}
\end{frame}

% ID: C08-I04
\begin{frame}{Bibliografía guiada}
\UASlideTag{RESPALDO}{UAIngNavy}
\footnotesize
\begin{itemize}
\item Kleppmann, \emph{Designing Data-Intensive Applications}, caps. 5, 8 y 9.
\item Lamport, \emph{Time, Clocks, and the Ordering of Events}.
\item Ongaro y Ousterhout, \emph{In Search of an Understandable Consensus Algorithm}.
\item Lamport, \emph{Paxos Made Simple}.
\item DeCandia et al., \emph{Dynamo}.
\item Corbett et al., \emph{Spanner}.
\end{itemize}
\end{frame}

\end{document}
